\begin{abstract}
Các hệ thống Data Pipeline hiện đại xử lý dữ liệu qua nhiều giai đoạn như thu thập, chuyển đổi, lưu trữ và phục vụ phân tích. Tuy nhiên, sau khi pipeline hoàn tất, nhiều nền tảng chưa có cơ chế độc lập để kiểm chứng rằng dữ liệu đã lưu vẫn khớp với trạng thái được sinh ra tại thời điểm thực thi. Nhật ký kiểm toán, checksum ở tầng ứng dụng và dashboard giám sát đều hữu ích, nhưng thường nằm trong cùng ranh giới tin cậy với dữ liệu hoặc chỉ phản ánh trạng thái khi pipeline đang chạy. Bài báo này trình bày một nghiên cứu đánh giá bán thực nghiệm đối với cơ chế \hashledger{} nhằm kiểm chứng tính toàn vẹn dữ liệu trong Data Pipeline. Cơ chế đề xuất ghi nhận bằng chứng ở từng stage, bao gồm số lượng bản ghi, schema hash, batch hash, block hash, liên kết previous hash và sự kiện lineage. Verification engine tính toán lại các bằng chứng sau xử lý và trả về trạng thái vi phạm toàn vẹn nếu có sai lệch. Nghiên cứu so sánh điều kiện sạch và điều kiện bị can thiệp, đánh giá khả năng truy vết đến first broken block và lineage event liên quan, đồng thời ước lượng overhead theo các mức record count khác nhau. Kết quả thực nghiệm trên nguyên mẫu Databricks cho thấy tất cả tamper scenarios được định nghĩa trước đều được phát hiện trong phạm vi kiểm thử, trong khi điều kiện sạch không tạo false positive trong dashboard export quan sát được. Overhead tăng theo record count, nhưng kết quả vẫn bị giới hạn bởi biến thiên môi trường cloud compute và số lần lặp benchmark còn hạn chế. Đóng góp của bài báo không phải là một hệ thống Blockchain đầy đủ, mà là đánh giá có kiểm soát một cơ chế hash-linked tamper-evident phục vụ bảo đảm tính toàn vẹn dữ liệu trong Data Pipeline.
\end{abstract}

\begin{IEEEkeywords}
Tính toàn vẹn dữ liệu, tamper-evident ledger, hash chain, data pipeline, bán thực nghiệm, an toàn thông tin, data lineage.
\end{IEEEkeywords}
